\documentclass{report}
\usepackage{amsmath,amssymb}
\setlength{\parindent}{0mm}
\setlength{\parskip}{1em}
\begin{document}
\begin{center}
\rule{6in}{1pt} \
{\large Ben, S. Southworth \\
{\bf A new look at interpolation in root-node smoothed aggregation}}

2805 Olson Dr \\ Boulder \\ CO 80303
\\
{\tt ben.southworth@colorado.edu}\\
Tom Manteuffel\end{center}

Currently, interpolation operators, P, in root-node smoothed aggregation
are formed through an energy minimization process over the columns of P.
Constraining P to exactly interpolate known near null-space candidate(s)
is a row-wise constraint, which together leads to an expensive, global
minimization process. In this talk, a new method to form P is proposed by
minimizing the $L^2$ error in interpolation. When combined with near
null-space constraints, the result is a local minimization process for
each row of P. A theoretical framework and two-level results are
discussed, motivating the next steps in achieving convergence rates of
current root-node at a much lower cost.


\end{document}
